\section{结论}
\label{sec:conclusion}

本文将叙事操纵不视为静态分类问题，而是视为通过语言、信念和网络化反应展开的动态过程。借助\textbf{NarrateSim}，我们可以在受控反事实条件下重播同一讨论，并询问一种策略如何变得有影响力、对谁有效、以及在何种结构条件下。更广泛的启示是：协调性叙事攻击是系统级威胁——说服性措辞固然重要，但受众信任特征、情绪状态、时机和局部网络位置同样重要。因此，保护在线社区不仅需要更好的检测器或在伤害扩散后进行更多事实核查，更需要将韧性设计到讨论环境本身中——结合早期认知准备、定向证据纠正和对协调放大的结构性限制。我们发布该框架和配置以支持后续在更多话题、语言和对抗设计上的研究。
